\section{ADR-0007: Dynamic Skip-Logic Scaffolding and Sheaf-Morphic Adaptive Trajectory Routing}
\label{sec:adr0007}

\begin{adrbox}{ADR-0007: Dynamic Skip-Logic Scaffolding}{Status: \statusbadge{ACCEPTED}}
\textbf{Date:} 2026-11-23 \\
\textbf{Context:} Runtime non-linear routing, interactive simulations, and cognitive intervention. \\
\textbf{Deciders:} Narrative Systems Team, Cognitive Engineering Group \\
\textbf{Status:} Accepted / Production Target
\end{adrbox}

\subsection{Context and Problem Statement}
Classical linear branching forces users down rigid, pre-authored dialogue trees or hardcoded failure loops \cite{nystrom2014game}. When a user encounters a cognitive barrier (such as conceptual difficulty with exponential growth or vector calculus), standard game engines either fail the user repeatedly or present simplistic text hints \cite{mourokh2024qari}. 

In contrast, production-grade skip-logic must dynamically detect cognitive stalls, route the user into scaffolded, interactive micro-simulations (e.g., visual PhET modules), and guarantee seamless return convergence to the primary curriculum objective without narrative discontinuity \cite{mourokh2024qari}.

\subsection{Decision Drivers}
\begin{enumerate}[label=\bfseries D\arabic*.]
    \item \textbf{Autonomous Scaffolding Detection:} Identify conceptual deficiencies automatically by decomposing state vector $\PsiVec$ into prime factor criteria.
    \item \textbf{Topological Convergence Guarantees:} Mathematically guarantee that all dynamic skip routes converge back to the primary trajectory (no orphan loops).
    \item \textbf{Interactive Micro-Module Injection:} Dynamically mount and unmount interactive simulation engines (e.g., PhET physics/math interactive simulations).
    \item \textbf{Narrative Preservation:} Maintain emotional immersion and narrative coherence across dynamic diversion branches.
\end{enumerate}

\subsection{Decision Outcome}
We implement a **Closed-Loop Sheaf-Morphic Skip-Logic Routing Engine**.

\subsubsection{Algorithmic Trigger Logic}
The skip-logic routing subsystem continuously monitors state vector $\PsiVec$ and the telemetry metrics $\RLG(t)$ and $\EFC(t)$:
\begin{enumerate}
    \item \textbf{Stall Detection:} A cognitive stall on prime concept $p_i$ is triggered when:
    \begin{equation}
    \alpha_{p_i}(t) < \tau_{\text{mastery}} \quad \text{and} \quad \frac{d}{dt}\alpha_{p_i}(t) \le \epsilon_{\text{stagnant}}
    \end{equation}
    \item \textbf{Morphic Diversion:} The engine activates the sheaf morphism:
    \begin{equation}
    \Phi_{\text{scaffold}}: \mathcal{U}_{\text{primary}} \longrightarrow \mathcal{U}_{\text{scaffold}}^{(p_i)}
    \end{equation}
    diverting user interaction from the primary scene into a targeted visual scaffolding module.
    \item \textbf{Scaffold Execution:} The Resource Subsystem mounts the interactive micro-simulation (e.g., a PhET vector addition simulation).
    \item \textbf{Convergence and Return:} As the learner interacts with the scaffolding module, $\alpha_{p_i}(t)$ increases. Once $\alpha_{p_i}(t) \ge \tau_{\text{return}}$, the inverse sheaf morphism $\Phi_{\text{return}}$ glues the learner back into the primary narrative trajectory.
\end{enumerate}

\begin{figure}[htbp]
\centering
\begin{tikzpicture}[scale=0.95, every node/.style={transform shape}, auto, >=Stealth]
    % Primary trajectory
    \node [draw, thick, rectangle, fill=primaryblue!15, rounded corners] (nodeA) {\textbf{Task Node A}\\Prime 2 Mastery};
    \node [draw, thick, rectangle, fill=primaryblue!15, rounded corners, right=4cm of nodeA] (nodeB) {\textbf{Task Node B}\\Primes $2 \cdot 3$ Composite};
    
    % Scaffold branch
    \node [draw, ultra thick, rectangle, fill=statusproposed!20, rounded corners, below=2cm of $(nodeA)!0.5!(nodeB)$] (scaffold) {
        \textbf{Scaffolded Sheaf: PhET Simulation Module}\\
        \footnotesize Prime $\ket{3}$ Targeted Scaffolding (Spatial Dynamics)\\
        Exit Criteria: $\alpha_3(t) \ge 0.90$
    };

    % Direct path (failed)
    \draw [->, thick, dashed, red] (nodeA) -- node[above, font=\footnotesize, text=red] {Direct path blocked: $\alpha_3 < 0.60$} (nodeB);
    
    % Divert arrow
    \draw [->, ultra thick, primaryblue] (nodeA) |- node[near start, left, font=\footnotesize] {Skip-Logic Trigger: $\Phi_{\text{scaffold}}$} (scaffold);
    
    % Return arrow
    \draw [->, ultra thick, statusaccepted] (scaffold) -| node[near end, right, font=\footnotesize] {Convergence: $\Phi_{\text{return}}$} (nodeB);
\end{tikzpicture}
\caption{Sheaf-Morphic Skip-Logic Trajectory Diversion and Convergence.}
\label{fig:skip_logic_routing}
\end{figure}

\subsection{Rust Routing State Machine}

\begin{lstlisting}[language=Rust, caption={Dynamic Sheaf-Morphic Skip-Logic Router}]
pub struct SheafMorphism {
    pub required_prime: u32,
    pub trigger_threshold: f32,
    pub return_threshold: f32,
    pub scaffold_bundle_id: &'static str,
}

pub struct SkipLogicRouter {
    morphisms: &'static [SheafMorphism],
    current_sheaf_id: u32,
}

impl SkipLogicRouter {
    pub fn evaluate_routing(&mut self, state: &PrimeStateVector, rlg: f32) -> Option<RoutingCommand> {
        for morphism in self.morphisms {
            let prime_idx = morphism.required_prime;
            if let Some(amp) = state.get_amplitude(prime_idx) {
                // Cognitive deficit detected; initiate dynamic skip logic
                if amp < morphism.trigger_threshold && rlg < 0.5 {
                    return Some(RoutingCommand::DivertToScaffold {
                        bundle_id: morphism.scaffold_bundle_id,
                        target_sheaf: prime_idx,
                    });
                }
            }
        }
        None
    }
}
\end{lstlisting}

\subsection{Consequences}
\subsubsection{Positive}
\begin{itemize}
    \item \textbf{Guaranteed Progression:} Players are never permanently stuck; targeted interactive scaffolds remediate specific concept deficits in real time.
    \item \textbf{Topological Safety:} Compile-time sheaf checks ensure all skip branches return to the main manifold.
    \item \textbf{Maximized RLG:} Dynamically minimizes wasted structural interventions, maximizing long-term learning gain.
\end{itemize}

\subsubsection{Negative and Mitigations}
\begin{itemize}
    \item \textbf{Asset Loading Latency:} Loading simulation bundles dynamically could cause hitching. \textbf{Mitigation:} Pre-cache high-probability scaffolding bundles in background threads based on predictive state trends.
\end{itemize}
